% !TEX program = xelatex
\documentclass{article}
\usepackage{graphicx} % Required for inserting images
\usepackage{amsthm, amsmath, amssymb, amsfonts} % general mathmode commands
\usepackage{bm} % for \bm
\usepackage{stmaryrd} % for \nnearrow
\usepackage{cancel} % for \cancel
\usepackage{breqn} % automatic line breaks in equations
\usepackage[margin=2.5cm]{geometry} % for page layout
\usepackage{cancel} % for \cancel
\usepackage{algorithm2e} % for algorithms
\usepackage{bbm} % for \mathbbm{1}
\usepackage{mathtools} % for \coloneq
\usepackage{hyperref}
\usepackage{nicefrac}
\usepackage{float}
\usepackage{polyglossia}
\usepackage{biblatex}
\usepackage{lipsum}
\usepackage{placeins}
\usepackage{booktabs}
\usepackage{xcolor,colortbl}
\usepackage{makecell}
\usepackage{tabularx}
\RestyleAlgo{ruled}

% Language settings
\setdefaultlanguage{hebrew}
\setotherlanguage{english}
\setmainfont[
  Path        = fonts/,
  Extension   = .ttf,
  UprightFont = *-Regular,
  BoldFont    = *-Bold,
  ItalicFont  = *-Medium
]{DavidLibre}

\title{סיכום אג''ם שגרה במד''א}
\date{}

\begin{document}
\maketitle
\section{אירוע רב נפגעים (אר''ן)}
אר''ן הוא אירוע בו מתקיימים אחד מהשניים:
\begin{enumerate}
    \item אירוע בו האמצעים ויכולות הטיפול בשטח האירוע אינם מספיקים לאורך זמן לטיפול בנפגעים.
    \item כמות הנפגעים וחומרת פציעתם עולות על יכולות הטיפול והפינוי בשגרה
\end{enumerate}
כמו כן, אר''ן מתחלק לשתי קטגוריות:
\begin{description}
    \item[קטגוריה ראשונה] אר''ן בו המערכת המרחבית יכולה להכיל את האירוע והנפגעים הדחופים מגיעים לבתי החולים עד שעתיים מזמן האירוע
    \item[קטגוריה שנייה] אר''ן בו לא ניתן לסיים את האירוע במהירות או באירועים בהם קיימת היערכות מוקדמת, כגון אירועים ציבוריים.
\end{description}
מבחינת מתן מענה לאר''ן, העיקרון המנחה הביסי הוא מתן מענה \textbf{ללא פגיעה בשגרה}.

\subsection{מענה רפואי לאר''ן}
בגלל אופי האירוע, העיקרון הבסיסי במענה הרפואי הוא שלא נותנים מענה רפואי בשטח מלבד מענה לשט''דים, ויתר המענה ינתן בשלב השניוני תחת סכימת \textenglish{PHTLS} ובבית החולים. בפרט, נגזר מכאן שהציוד המורד מהאמבולנס הוא כלי נשיאה (מיטה, כיסא, לוח גב) וכלי חבישה (תחבושות אישיות, חסמי עורקים), וכן תגי מיון עבור הנפגעים. במהלך המעבר בין החולים, החובשים יגדירו וסימנו האם המטופלים במצב דחוף או לא דחוף (לפי פגיעה ב-\textenglish{ABC}), והפראמדיקים/רופאים יבצעו מיון שניוני.
\subsection{מענה אג''מי לאר''ן}
איש הצוות הבכיר באמבולנס הראשון שמגיע לאירוע הוא \textbf{פיקוד 10} כלומר מפקד האירוע עד קבלת ההחלטה להעביר את הפיקוד. לפיכך, התנהלותו הראשונית מתווה את המשך הפעילות באר''ן.

בחירת חניית האמבולנס הראשון מושתתת על מספר עקרונות:
\begin{enumerate}
    \item בטיחות
    \item קרבה למקום האירוע
    \item מיקום בולט עבור נקודת ריכוז
    \item הימנעות מחסימת צירים
    \item מקום לריכוז נפגעים
    \item מקום לאמבולנסים נוספים
\end{enumerate}
הדיווח הראשוני יהיה:
\begin{enumerate}
    \item מיקום מדויק של האירוע
    \item אופי האירוע/סיכונים בולטים (תאונת דרכים, פיצוץ, ירי, שריפה, קריסה)
    \item הערכת מספר הנפגעים (בודדים, עשרות, מאות)
    \item צירי הגעה לאירוע/צירים חסומים, נקודת כינוס ומיקום חניית האמבולנס
    \item הכרזה על פיקוד 10
\end{enumerate}

על מנת לבצע את המיון באופן יעיל, פיקוד 10 יכרוז לנפגעים מהלכים וירכז את כלל הנפגעים בנקודת ריכוז, ינהל קו פתוח מול המוקד המפעיל, יחלק את השטח לגזרות ברורות וימנה אחראים, יפעיל את הצוותים הרפואיים, ידאג לפיזור וויסות המטופלים לבתי החולים וינהל סריקות וסריקות חוזרות. כמו כן ינהל את תמונת המצב וישלוט במשאבים, וימנה מפקדי משנה.

\subsection{מפקדי משנה באר''ן}
\subsubsection{פינוי 10}
ימוקם בצוואר בקבוק פינויי, יוודא הימצאות ציוד תמאים על האמבולנס, ירכז את תמונת המצב של הפינויים שאינם של מד''א ויסייע למוקד המפעי לביצירת תמונת מצב פינויים.
\subsubsection{רפואה 10}
ינהל את מיון הנפגעים לפי הדחיפות, ינחה על ביצוע טיפולים מצילי חיים בנפגעים, יוודא קדימויות, יווסת צוות רפואי מתקדם, יוודא קביעת או הכרזת מוות.
\subsubsection{חניון 10}
ירכז את רכבי ההצלה בנקודה שהוגדרה, יוודא שמשאירים מפתח בפנים וחלון נהג פצוח, יוודא את מיקום החניה, יפעיל את הצוותים, ציוד ורכבי ההצלה.
\subsubsection{נציג מד''א בביה''ח}
תפקיד מוגדר מראש שמתייצב בביה''ח בעת אר''ן, יבצע רישום הגעת הנפגעים ויעדכן את בית החולים לגבי נפגעים שבדרך, מצבם וסיום פינוי ואת הממל''מ על יכולות קליטה ותמונת מצב הנפגעים, וכן יחזיר ציוד של מד''א לתחנות בסיום האירוע.

\subsection{דגשים}
\begin{enumerate}
    \item באירועי פיצוץ אסור לפנות חפצים מהשטח ויש להודיע לחבלן על נקודת ריכוז הנפגעים.
    \item באירועי קריסה יש לחשוש לפציעות מעיכה
\end{enumerate}

\section{חומרים מסוכנים (חומ''ס)}
\subsection{ניהול החומרים}
חומ''ס הוא כל חומר שעלול לגרום לנזק לאדם או לסביבה. ועדת המומחים של האו''ם להובלת חומרים מסוכנים (\textenglish{UN TDG Committee}) מנהלת ומתחזקת רשימה כל כלל החומרים המסוכנים, המזוהים ע''י ארבע ספרות, למשל \textenglish{1203} (דלק) וכד'. החומרים מחולקים ל-\textenglish{8} מחלקות לפי מהות החומר, כגון חומר נפץ, חומר מחמצן, חומר רעיל וכד'. במד''א ישנו ``הספר הכתום'' שמכיל הדרכות רפואיות לטיפול בחומרים מסוכנים רבים.
\subsection{אירועי חומ''ס}
 אירוע חומ''ס הוא התרחשות בלתי מבוקרת או תאונה שמעורב בה חומ''ס הגורמת או עלולה לגרום סיכון לאדם ולסביבה, לרבות שפך, דליפה, פיזור, פיצוץ, התאיידות, דליקה. ככלל נזהה אירוע חומ''ס ע''י ריח חריף בזירה, מראות שריפה, מספר או''ם, עננים נמוכים וחריגים, או סימני חשיפה פיזיים של המטופלים. אירוע החומ''ס השכיחים ביותר שמד''א מוזנק אליהם נגרמים עקב טעויות אנוש או רשלנות הקשורות לגז הבישול בצוברים או במיכלים הביתיים.

\subsection{בהגיענו לאירוע}
ראשית, נקבע מרחק ביטחון ראשוני (מב''ר) לפי סוג אירוע החומ''ס: אירוע בשטח סגור - 100 מטרים מפתחי המבנה, אירוע בשטח פתוח - 200 מטר ממוקד האירוע, אירוע גז טבעי/מתאן - 300 מטר ממוקד האירוע.
המב''ר מהווה קו גבול ראשוני שמסומן בטאבלט ע''י המוקד המרחבי שאסור להיכנס אליו. לאחר הערכת סיוכנים ע''י כב''א, במה''ר יתבטל ויחלף בחלוקה לשלושה אזורים: קר, פושר וחם. אזור חם הוא אזור בו רמת הזיהום והריכוז של החומרים הרעילים גבוהה. מד''א יכולים לפעול בשטח זה רק לאחר אישור הגורמים המוסמכים. באזור פושר רמת הרעילות נמוכה יחסית אך חשיפה אליה עלולה לגרום לנזק בריאותי בחומרה קלה עד בינונית. בשטח זה צוותי מד''א יוכלו לפעול עם ציוד מגן תואם ולזמן קצר. אזור בו רמת הזיהום והריכוז של החומרים
המעורבים בזירת האירוע נמוכה וניתן לפעול באופן חופשי ללא מיגון.
\subsection{מיגון}
ערכת המיגון האישי (C+) שלנו מורכבת מכמה חלקים:
\begin{enumerate}
    \item חליפת הגנה מפני התזה (סי''ל)
    \item מסיכה אקטיבית + מסנן ABEK2P3
    \item כפפות
    \item ערדליים
    \item טלית זיהוי
\end{enumerate}
\subsection{שטיפה}
על מנת לטפל בנפגעים, חשוב לשטוף אותם כדי להסיר את שאריות החומ''ס. אם החומר זוהה, נפעל לפי ההנחיות הפרטניות. אם החומר לא זוהה, נפעל לפי השיטה הבאה:
\begin{description}
    \item[גז] הפשטה בלבד
    \item[נוזל] הפשטה ושטיפה במים זורמים
    \item[אבקה] הפשטת הנפגעים והסרת האבקה מהנפגע עד כמה שניתן
\end{description}
כמו כן, כלל צוותי החירום היוצאים מהאזור המזוהם חייבים להישטף טרם הסרת חליפת המיגון שלהם.
\subsection{כניסה לאזור מזוהם}
ניכנס לאזור מזוהם פושר/חם רק לצורך הצלת חיים ורק במידה שיש נפגעים המראים סימני חיים. ניכנס עם שני אנשי צוות מטעמנו המצוידים בציוד מיגון אישי הקיים ברכבי ההצלה בליווי צוות כב''ה עד לאזור בו הנפגעים מראים סימני חיים. המטרה שלנו היא הפסקת חשיפת הנפגעים לחומר המסוכן וביצוע פעולות מצילות חיים בנפגעים לכודים או במקביל לחילוצים. הכניסה תימשך עשרים דקות לכל היותר.
\section{קודים מבצעיים}
למד''א סרגל אחיד של קודים מספריים בין 1 עד 9000 שמשמש עבור כלל הקודים המבצעיים, קודים לרכבים, קודים לתפקידים, קודים לאירועים וכו'. הקודים משמשים לתקשורת בקשר המבצעי. לכל טווח של קודים יש משמעות, למשל קודים בין 11 ל-399 הם בדר''כ קודים לנט''נים, קודים בין 800-2500 הם בדר''כ קודים לאמבולנסים וכו'. קודי רכב הם חלק מהקודים המבצעיים שהם מזהים רכב באופן ייחודי. הם מסומנים על הרכב בכל צדדיו. בשיח בקשר משתמשים במרחב-קוד מבצעי, למשל שרון-2005 או ירקון-33.
\section{אירוע 100}
מד''א, בהיותו ארגון ההצלה הלאומי של מדינת ישראל נותן מענה רפואי לקריאות המתקבלות מהיחידה לאבטחת אישים חשובים מאוד (אח''מ) במשרד ראש הממשלה. מד''א 100 הוא הקוד המבצעי של היחידה לאבטחת אישים בעת הזעקת סיוע מד''א, ומד''א 99 הוא הקוד המבצעי לפאראמדיק מד''א בצוות היחידה לאבטחת אישים. בעת אירוע 100, צוות מד''א ישמע להוראותיו של מפקד האבטחה מהיחידה לאבטחת אישים. בעל הסמכות הרפואית העליונה הוא רופא היחידה לאבטחת אישים. אם זה אינו נוכח באירוע, יהיה בעל הסמכות פאראמדיק היחידה, למעט אם נוכח רופא בצוות מד''א
שהגיע למקום. ככלל, יפונה האח''מ למרכז טראומה על אזורי.

\subsection{אירוע מתוכנן}
לאירוע ידוע מראש של אח''מ שצריך אבטחה רפואית, יוזמן מראש אט''ן שנקרא אט''ן אח''מ. פרטי העובדים/מתנדבים על האט''ן יועברו מראש ליחידת האח''מ ואסור להחליפם ללא אישור. צוות אט''ן אח''מ יחבור למפקד האבטחה/רופא היחידה לתדרוך על מקום החנייה, ציר הפינוי ויעד הפינוי ויהיה כפוף להנחיותיו במהלך האירוע. צוות אט''ן אח''מ ישאר ברכב בהאזנה לגל המרחבי, ואם סופק לו מכשיר קשר של היחידה יאזין לו במקביל. במקרה של בקשה אקראית לעזרה רפואית (לא לאישיות המאובטחת), צוות אט''ן אח''מ יתחיל בטיפול רפואי אך יבקש
מהמוקד המרחבי לשגר צוות אחר להמשך הטיפול והפינוי ויתאם את נושא הטיפול עם מפקד האבטחה.

\subsection{קריאה לאירוע 100}
בעת אירוע 100, מאבטח או גורם אחר ביחידה יזעיק את מד''א באמצעות טלפון ויזדהה בכינוי מד''א 100. פאראמדיק היחידה עשוי להזדהות בהתקשרות בכינוי מד''א 99. ההודעה תימסר למד''א בלשון גלויה ותכלול מס' טלפון (להתקשרות חוזרת), תיאור מקרה, מס' נפגעים, אופי הפגיעה, מיקום האירוע או אתר לחבירה. לאחר שנקבע יעד הפינוי יעדכן המוקד המרחבי את ביה''ח הקולט על העברת המטופלים, מצבם וזמן הגעה משוער.

\section{אירוע 700}
משרד הביטחון (משהב''ט) הוא המשרד הממונה על הפעלת הכוח במדינת ישראל. בין היתר משרד הביטחון אמון על הצבא ועל יחידות נוספות. בתור ארגון ההצלה הלאומי, מד''א נקרא לתת סיוע רפואי למשהב''ט בשני מתארים עיקריים:
\begin{enumerate}
    \item סיוע רפואי ליחידות משהב''ט הנמצאות בפעילות
    \item סיוע רפואי באירוע המתרחש במתקנים בשייכים למשהב''ט
\end{enumerate}
באירועים כאלה, מד''א 700 הוא כינוי ההפעלה של משהב''ט

\subsection{סיוע רפואי ליחידות משהב''ט}
בעת הזעקת סיוע רפואי, מוסר הדיווח יזדהה כמד''א 700 ויעביר למוקד את מיקום האירוע/נקודת חבירה, מספר הנפגעים ומצבם, אופי האירוע וגורמי סיכון סביבתיים ומספר טלפון (להתקשרות חזרה). בפינוי המטופל יעדכן צוות מד''א את הגורם האחראי בשטח ממשהב''ט ביעד הפינוי. קציני המשמרת במוקד המרחבי והארצי ידווחו מיידית על הפעלת מד''א בסיוע למשהב''ט.
\subsection{סיוע מד''א באירוע במתקני משהב''ט}
ברחבי המדינה פרוסים מתקנים של משהב''ט. מתקנים אלה רגישים מבחינה בטחונית. אם נדרש סיו עמד''א באחד ממתקני משהב''ט, עשוי המתקן להתקשר למוקד הארצי או ישירות למוקד המרחבי שבגזרתו התרחש האירוע. המוקד המרחבי ישגר בנסיעה דחופה את הצוות הזמין הקרוב לאירוע (בעדיפות אט''ן) וכוחות נוספים בהתאם לצורך ויגדיר במערכת השו''ב אירוע מסווג קוד 62, וכן ידווח על האירוע למוקד הארצי.
\section{הסקה}
הפעלת מסוק לצורך טיפול ופינוי, בדגש על אזורים מרוחקים ופריפריאליים, משפרת את הנגישות למטופל ל-ALS ועשויה לקצר את הזמן הנדרש לפינוי הנפגע מאזורים אלו. מד''א מפעיל שני מסוקים אזרחיים (צפוני ממנחת פוריה ודרומי ממנחת דבירה) וכן רשאי להסתייע במערך המסוקים של צה''ל.
\subsection{שיטת הפעלה של מסוק אזרחי}
במטוס אזרחי צוות של טייס, פראמדיק ראש צוות ופראמדיק מוטס.
הסמכות והאחריות של המסוקים היא של המוקד הארצי. הסמכות לבקש סיוע של מסוק או לבקש ביטול שיגור המסוק ממוקד מרחבי/ארצי היא של איש הצוות הבכיר הנוכח באירוע. כמו כן, המוקד המרחבי יכול לבקש מהמוקד הארצי הזנקת המסוק.

במסוק, הטייס הוא מפקד המסוק והמשימה והפראמדיק ראש הצוות הוא האחראי על הרפואה.

ככלל, מטוס יוזנק לטובת טיפול \textenglish{ALS} כאשר משך הזמן שידרש לצוות המסוק להגיע אל החולה/נפגע, לטפל ולפנות לביה''ח הוא הקצר ביותר, לצורך שינוע דחוף של מנות דם או איברים להשתלה, וכן במקרים אחרים בהוראת מנכ''ל מד''א או ראש אג''ם. כמו כן, לצורך העברה לא דחופה יפנה המוקד של המבקש לפקמ''צית חב' המסוקים להזמנת השירות, זאת באיושר סמנכ''ל רפואה ובהיעדרו כונן אגף רפואה.

יש לשים לב שמשך הזמן הנדרש למסוק כדי להמריא עומד על כ-10 דקות מרגע ההודעה לצוות, וכן במצב מיוחד של ``קיצור כוננות'' זמן ההמראה מתקצר לכ-5 דקות. לבסוף זמן החסרת המסוק לכשירות מבצעית בתום משימתו ולאחר נחיתתו היא כחצי שעה.

יחד עם זאת, מטופל לא יפנה במסוק כאשר יש מוות ודאי, קביעת מוות ע''י רופא או הפסקת החייאה ע''י פראמדיק, נפגע חומ''ס שעל גופו שארית חומר מסוכן ואי אפשר לטהרו או מטופל משתולל.
\subsection{שיטת הפעלה של מסוק צבאי}
במסוק צבאי צוות של שני טייסים, רופא ומספר חובשים. בדומה למסוק אזרחי, הסמכות לבקש מסוק צבאי היא של איש הצוות הבכיר באירוע, שבתורו פועל מול קשל''ט מסוקים. יש לשים לב שזמן ההכנה של מסוק צבאי הוא ארוך - יציאה לדרך תוך 15 דקות מרגע הבקשה בניגוד לעשר דקות במסוק אזרחי או חמש דקות במסוק אזרחי עם קיצור כוננות.
\section{עבודת המוקד}
מוקד מרחבי וארצי
\section{שו''ב ושליטה}
\section{אירועי בני ערובה}
אירוע בני ערובה הוא אירוע בו גורמים עוינים לוקחים בשבי קבוצה של בני ערובה ומאיימים לפגוע בהם. לפיכך, זה אירוע מורכב גם ברמה הבטחונית וגם ברמה הרפואית. על מנת לטפל באירוע בני ערובה ברמה הבטחונית, מפעילים כל אחד מהגופים הבאים: צה''ל, מ''י ושב''ס יחידת השתלטות משלו שתפקידה לפעול במקרה של אירוע בני ערובה ולנטרל את האיום הבטחוני תוך שמירה על חייהם ושלומם של בני הערובה ככל האפשר. בין חברי היחידה נמנה לפחות פאראמדיק או רופא אחד. יחד עם זאת, לאור המורכבות הרפואית, צוותי מד''א ידרשו
לסייע באירוע בהיבטו הרפואי.

ביחס לאירוע בני ערובה, המרחב מחולק למעגל פנימי - תחום פעילותה של יחידת ההשתלטות ומעגל חיצוני - כל מה שאינו נכלל במעגל הפנימי. ביניהם נמצאת נקודת מיון ראשונית אליה מפונים הנפגעים, ורופא/פאראמדיק מתאם מטעם יחידת ההשתלטות אמון עליה. כמו כן באירוע שצה''ל מטפל בו, ישנה נקודת רפואה מרכזית שממוקמת במעגל החיצוני ומאוישת ע''י רופא ומספר חובשים מצה''ל.

באירוע בני ערובה, הגורם הפוקד הוא צה''ל, מ''י או שב''ס, והאחריות למענה הרפואי במעגל הפנימי מוטלת על יחידת ההשתלטות מטעם הגורם הפוקד. האחריות הרפואית בנוקת המיון הראשונית מוטלת על הרופא/הפאראמדיק המתאם.

במעגל החיצוני, באירועים של שב''ס ומ''י האחריות הרפואית מוטלת על מד''א, ובאירוע של צה''ל היא מוטלת על צה''ל, שמקים נקודת רפואה מרכזית.


\section{אירועים בנתב''ג}
נתב''ג מהווה פוטנציאל איום להתרחשות אר''ן על רקע תאונת מטוס או פעילות חבלנית עוינת. בגלל רגישות האתר, מד''א ערוך באופן מיוחד לאירוע חירום בו. בפרט, מד''א מציב עתודת רכבי חירום בסמוך לשער ברוש.

בהפעלה, ישנו סרגל של מצבי הפעלה על פי חומרת התקלה: כוננות \(\leftarrow\) מצב חירום 2 \(\leftarrow\) מצב חירום 3 \(\leftarrow\) ניצוץ אנושי וכן על כמות הנפגעים מטוס קל \(\leftarrow\) מטוס צהוב \(\leftarrow\) מטוס לבן \(\leftarrow\) מטוס ירוק \(\leftarrow\) מטוס אדום, כאשר האמצעים המופעלים מסלימים על פי העלייה בסרגלים.

הפעילות בנתב''ג מתבצעת בדבוקות שהן טור של שני נט''נים, 3-5 אמבולנסים ותאר''ן שמתניידים כיחידה אחת. 
\section{השימוש בטאבלט}
בכל אמבולנס מוצב טאבלט עם אפליקציה ייעודית שמשמש לשליטה ובבקרה באמבולנס ובאירועים השונים. לאפליקציה יש מספר מודולים:
\subsection{טיפול באמבולנס בתחילת המשמרת}
ראשית יש למלא טופס בדיקת אמבולנס שמונה את כל הפריטים שצריכים להיות באמבולנס לפי התקינה, ולכתוב כמה יש בפועל. כמו כן יש לציין נזקי פח באמבולנס.

\subsection{טיפול במאורע}
מאורע יתקבל בדחיפה לאפליקציה ויכיל את ויזת הנסיעה לפי הידוע למוקד ומיקום, עם אפשרות להכוונה ב-waze. לכל מאורע יש לפחות מטופל אחד, לו צריך למלא דוח מטופל בלחיצה.

\subsection{דוח מטופל}
\subsubsection{פרטי המטופל}
פרטי המטופל הם ת''ז, קופ''ח, שם, תאריך לידה ופרטים ליצירת קשר. לאחר הקשת ת''ז ניתן ללחוץ על הכפתור ''חפש'' והמערכת תמלא את יתר הפרטים במידת האפשר באופן אוטומטי.
\subsubsection{אנמנזה}
במסך האנמנזה יש שתי לשוניות: תלונה עיקרית ומחלה נוכחית:
\begin{description}
    \item[תלונה עיקרית (\textenglish{chief complaint})] בגין מה נקראנו לאירוע - אמור להיות קצר מאוד, משפט אחד ומטה.
    \item[מחלה נוכחית (\textenglish{present illness})] גוף האנמנזה. מורכב מארבעה חלקים שנכתבים כפסקה אחת: 
        \begin{description}
            \item[גיל ורקע רפואי] במילה אחת או שתיים
            \item[בהגיענו למקום] תיאור הסצינה, מה ראינו, כיצד המטופל נמצא, באיזו תנוחה, בנוכחות מי הוא נמצא.
            \item[לדברי המטופל] מה התסמינים שהמטופל מעיד עליהם, מה הוא מרגיש, מה הוא חווה.
            \item[בבדיקתנו] תוצאות הבדיקה באופן כללי (ללא מספרים) לפי הסדר: הכרה (צלול? מגיב?), נשימה (קושי נשימתי?), דם (לחץ דם תקין? דופק תקין? סדיר? כיצד נראה העור? מילוי קפילרי?) וכן תוספות נוספות לפי המקרה (האם יש סיפור של אבדן הכרה? האם יש חבלות ודימומים? לרשום אם אין חבלות ודימומים נוספים הנראים לעין).
            \item[פונה] להמשך טיפול ובירור רפואי וכד'. לציין לאן פונה או אם לא פונה.
        \end{description}
\end{description}

\subsubsection{מדדים}
צריך לבדוק את המטופל בקבלתו, וכן כל רבע שעה לאחר מכן. הבדיקות המבוצעות בכל סבב בן נשימות, דופק ולחץ דם וכן בדיקות נוספות אם הן רלוונטיות לפי הקליניקה של המטופל.
מבצעים וממלאים עמודת בדיקות כל רבע שעה כאשר המטופל ברשותנו (לפני ההגעה ליעד הפינוי), כאשר הבדיקות הקבועות הן הכרה, נשימה, לחץ דם ודופק (סוכר רק אם רלוונטי). כמו יש יש בדיקות בסיום שלוש דקות לפני הגעה ליעד הפינוי.

\section{נספח}
\subsection{רשימת טווחי קודים מבצעיים שמשמשים לתקשורת מבצעית}
\begin{table}[H]
    \caption{רשימת טווחי קודים מבצעיים שמשמשים לתקשורת מבצעית}
    \center
    \begin{tabular}{cc}
        \toprule
           טווח קודים & משמעות                                    \\
\midrule
                   1 & מנהל המרחב                                \\
                 2\textendash 7 & סגנים ועוזרים של מנהל המרחב                  \\
                   8 & סופרוייזר                                    \\
                   9 & קצב''ת (קצין בטיחות בתעבורה)                 \\
                  10 & מפקד באר''ן                                 \\
 11–399 למעט 99,100 & נט''נים                                     \\
                  99 & פראמדיק ביחידה לאבטחת אישים                \\
                 100 & היחידה לאבטחת אישים                        \\
             400–500 & נט''נים ממ''י                                 \\
             550–580 & חפ''קים                                     \\
             591–595 & מסוקים                                     \\
             600–699 & אמבולנסים 4*4                              \\
                 700 & אירוע משהב''ט                               \\
             702–799 & אמבולנסים ממוגנים                            \\
 800–2349 למעט 1000 & אמבולנסים רגילים                             \\
                1000 & אירוע אלימות כלפי צוות מד''א                   \\
           2350–2399 & טנ''צ (טיפול נמרץ צבאי)                      \\
           2400–2500 & אמבולנסים רגילים                             \\
           2600–2649 & תאר''נים                                    \\
           2650–2749 & טנדרים/ג'יפים ממ''י                           \\
           2750–2799 & טרקטורונים/טומקר                            \\
           2800–2899 & ג'יפים                                      \\
           2900–2949 & רכבי שירות                                  \\
           2950–2999 & רכבים אחרים – אוטובוס, אמבולנס המשאלות, סגווי \\
           3000–3400 & אמבולנסים יישוביים                            \\
           4001–4199 & ניידות דם                                   \\
           4200–4250 & ניידות דם ממ''י                              \\
           5000–5999 & אופנועים                                    \\
           6000–6200 & אופנועים של עמותות אחרות                     \\
           6201–6800 & רכב חשמלי                                  \\
           7500–9000 & אופניים                                    \\
        \bottomrule
\end{tabular}
\end{table}
\subsection{תשובות ללומדת מענה לאירוע רב נפגעים}
\begin{itemize}
\item הלומדה הראשונה בלומדות אג''ם (אר''ן) ארוכה מאוד. להלן מפתח תשובות למבחן מסכם.
\item כל השאלות עד המבחן המסכם לא נספרות לציון אז אפשר להעביר אותן.
\item לפני המבחן המסכם יהיה כיתוב גדול באדום של ''מבחן מסכם''.
\item במבחן יש שמונה שאלות ורק הן נקבעות לציון.
\item השאלות:
    \begin{enumerate}
        \item מהי ההגדרה של אירוע רב נפגעים? (יותר מתשובה אחת נכונה)

            תשובות: אירוע בו האמצעים ויכולת הטיפול אינם מספיקים לאורך זמן לטיפול בנפגעים; אירוע בו כמות הנפגעים וחומרת פציעתם גדולות ביחס ליכולות הטיפול והפינוי בשגרה.
        \item מי רשאי להכריז על אירוע רב נפגעים (יותר מתשובה אחת נכונה)

            תשובות: איש הצוות הבכיר באמבולנס מד''א הראשון; מוקד מד''א המקבל את האירוע
        \item מהן הפעולות הנדרשות מהצוות שבדרכו לאר''ן (יותר מתשובה אחת נכונה)

            תשובות: שימוש בפנקס כיס/סד''פ אר''ן הנמצאים באמבולנס לתדרוך הצוות; הכנת אמצעי קשר (טלפון סלולארי, POC); התמגנות, קבלת הנחיות בשיחות בהתאם לתרחיש ולבישת וסטים כאמצעי זיהוי; חלוקת משימות בין אנשי הצוות, לרבות ציוד המורד מהרכב.
        \item ישנה חשיבות רבה לבחירת מקום החניה של האמבולנס הראשון. אלו מהתשובות הבאות \textbf{לא נכונה} לגביה?

            תשובות: מקום חניית האמבולנס נקבע ונמסר ע''י המוקד המבצעי

       \item אתה נהג אמבולנס, המגיע כאחד מהצוות הראשונים (לא הראשון) לזירת אר''ן. איזה ציוד תוציא ממנו?

            תשובות: תגי מיון וטיפול, חסמי עורקים, אמצעי חבישה, מלע''כ, מיטת אמבולנס, אלונקה ולוחות גב.

            הסבר: לאר''ן לא מוציאים ציוד החייאה (דפיברילטור, אמבו) אבל מוציאים ציוד לשינוע מטופלים (אלונקה, לוחות גב).
        \item מהן הפעולות הרפואיות הדחופות ביותר להן נדרש איש צוות מד''א בעת טיפול בזירת פיצוץ?

            תשובות: ביצוע פעולות מצילות חיים וחילוץ הנפגעים תוך הרחקתם מסכנה.
        \item לפנינו מספר משפטים הנוגעים לפעילות צוותי מד''א בזירת פיצוץ. מהם המשפטים הנכונים ביותר? (יותר מתשובה אחת נכונה).

            תשובות: ''פיקוד 10'' יחבור לחבלן לתיאום נקודת ריכוז ויעדכן הנחיות בטיחות תוך הרחקת הנפגעים; נחנה את האמבולנס כ-50 מטר מהזירה ונתמגן באפוד מגן וקסדה.
        \item לפנינו משפטים בנושא תג מיון ופינוי נפגעים. מהו המשפט ש\textbf{אינו נכון}?

            תשובות: נשלים את כל הרישום בתג המיון, לפני פינוי הנפגע מהזירה.
    \end{enumerate}
\end{itemize}
\subsection{קודים חשובים בקשר}
\begin{center}
    \begin{tabularx}{\textwidth}{XX}
        \toprule
        קוד & משמעות \\
        \midrule
        99 & פראמדיק יחידת אח''מ \\
        \midrule
        100 & יחידת אח''מ \\
        \midrule
        700 & אירוע של יחידת משהב''ט או במתקן של משהב''ט \\
        \midrule
        צבא שמיים & כינוי לאירוע בני ערובה המוחזקים בידי מחבלים בישראל \\
        \midrule
        מיניסטר חוף &  כינוי לאירוע בני ערובה המוחזקים בידי מחבלים בכלי שיט \\
        \midrule
        כור היתוך & כינוי לאירוע בני ערובה המוחזקים בידי עצירים/אסירים במתקני כליאה/בתי מעצר \\
        \midrule
        טוהר לב & כינוי לאירוע בני ערובה המוחזקים בידי מחבלים במטוס \\
        \midrule
        מעגל פנימי & השטח בו פועלת יחידת ההשתלטות באירוע בני ערובה \\
        \midrule
        מעגל חיצוני & השטח החיצוני לפעילות יחידות ההשתלטות, בו פועלים בין היתר כוחות הרפואה \\
        \midrule
        נקודת מיון ראשונית & הנקודה המשיקה בין המעגל הפנימי לחיצוני אליו מפונים הנפגעים \\
        \midrule
        נקודת רפואה מרכזית & נקודה המופעלת ע''י צה''ל באירוע בני ערובה בהם הוא פוקד במעגל החיצוני \\
        \midrule
        התראה & מצב כוננות באירוע בני ערובה בו לא משוגר כוח מד''א \\
        \midrule
        מגש & כינוי לאירוע כתוצאה מפעילות חבלנית עוינת בשטח הקרקעי הציבורי סטרילי בנתב''ג \\
        \midrule
        מטוס קל & אירוע בנתב''ג שצפויים בו עד 10 נפגעים \\
        \midrule
        מטוס צהוב & אירוע בנתב''ג שצפויים בו עד 30 נפגעים \\
        \midrule
        מטוס לבן & אירוע בנתב''ג שצפויים בו עד 100 נפגעים \\
        \midrule
        מטוס ירוק & אירוע בנתב''ג שצפויים בו עד 200 נפגעים \\
        \midrule
        מטוס אדום & אירוע בנתב''ג שצפויים בו מעל 200 נפגעים \\
        \midrule
        כוננות & באירוע בנתב''ג, מוכרז כאשר מדווחת תקלה במטוס שעלולה לסכן את שלום נוסעיו. המענה לאירוע ינתן ע''י כוחות פנימיים של רשות שדות התעופה בלבד \\
        \midrule
        מצב חירום 2 & הכרזה על תקלה במטוס שיש בה בכדי לסכן את שלום נוסעיו עקב נחיתה. המענה ינתן בשילוב של כוחות הצלה פנימיים וחיצוניים \\
        \midrule
        מצב חירום 3 & מוכרז בהתרחשות תאונת מטוס או כאשר הנתונים מצביעים על תאונת מטוס בלתי נמנעת. במצב זה ניתן מענה ע''י כוחות חיצוניים ופנימיים \\
        \midrule
        ניצוץ אנושי & מוכרז במצב של התרסקות מטוס ואר''ן כתוצאה ממנו \\
        \midrule
        אר''ן & אירוע רב נפגעים \\
        \midrule
        אט''ה & אירוע טוקסיקולוגי המוני \\
        \midrule
        חומ''ס & חומר מסוכן \\
        \midrule
        פיקוד 10 & כינוי למפקד אירוע מטעם מד''א באר''ן \\
        \midrule
        מד''א 1000 & כינוי לקריאה לסיוע משטרתי לאור אלימות \\
        \midrule
        מב''ר & מרחק בטיחות ראשוני, המרחק הראשוני שיש לשמור ממוקד אירוע חומ''ס לפני ניתוח של כב''א \\
        \bottomrule
    \end{tabularx}
\end{center}
\end{document}

